% =============================================================================
% Appendix B — Reproducibility
% =============================================================================
\section{Reproducing Everything in This Book}\label{app:repro}

Every number, figure and theorem in this book is produced by a named command.
This appendix lists them, in the order you would run them.

\subsection*{1. Retrieve the data}

\begin{lstlisting}[language=bash,style=lean,morekeywords={python}]
python scripts/textbook/retrieve_genomic_data.py
\end{lstlisting}

Fetches gene records from Ensembl REST, variant counts from NCBI E-utilities,
and protein and domain records from UniProt and InterPro. Writes
\file{textbook/data/genomic_records.json} with the source URL and a UTC
timestamp beside every record.

The script verifies its own output before writing: every gene must be on GRCh38,
every domain must lie inside its protein, no ClinVar category may exceed the
total, and residue $286$ must fall inside the retrieved boundaries of
\texttt{IPR006133}. If any check fails it exits non-zero and writes nothing, so
the book cannot be built on records that contradict its claims. Use
\texttt{-{}-offline} to re-run the checks against the cache without network
access.

\subsection*{2. Build the figures and numerical results}

\begin{lstlisting}[language=bash,style=lean,morekeywords={python}]
python scripts/textbook/make_figures.py            # all figures
python scripts/textbook/make_figures.py --only 3 5 # selected sections
\end{lstlisting}

Writes PDFs to \file{textbook/figures/} and every quantity quoted in the text to
\file{textbook/data/figure_values.json}, so prose and plots cannot drift apart.
Section~8 additionally caches the \texttt{POLE} coding sequence it fetches.

\subsection*{2b. Run the two experiments of Section~\ref{sec:gaps}}

\begin{lstlisting}[language=bash,style=lean,morekeywords={python}]
python scripts/textbook/tda_sequence_null.py --shuffles 200
python scripts/textbook/verify_yeast_human_correspondence.py
\end{lstlisting}

The first runs the shuffle null of \S\ref{sub:tdanull}: three nested nulls
against the $H_1$ total persistence of each coding sequence, with an i.i.d.
negative control and a tandem-repeat positive control. The second runs the
full-length alignment of \S\ref{sub:ortho} under four scoring schemes, with a
self-alignment positive control and a composition-matched-shuffle negative
control.

Both \textbf{assert} their controls and exit non-zero without writing if either
fails, so a broken test cannot quietly emit a $p$-value. The first also verifies
that each shuffle preserves the invariant it claims --- base counts,
dinucleotide counts, codon usage --- before any null is drawn. Results land in
\file{data/tda_sequence_null.json} and
\file{data/yeast_human_correspondence.json}, including the individual shuffle
draws so that Figure~\ref{fig:tdanull} plots the real distribution rather than a
Gaussian fitted to its first two moments.

The null is the slow step: roughly fifteen minutes on eight cores at $200$
shuffles. Use \texttt{-{}-shuffles 20} for a smoke test, remembering that the
attainable $p$ floor is then $2/21 = 0.095$ and the positive control cannot
reach $0.05$.

\subsection*{2c. Run the four experiments of Section~\ref{sec:reticulata}}

\begin{lstlisting}[language=bash,style=lean,morekeywords={python,curl,cargo,cd}]
cd rust/reticulation && cargo build --release
./target/release/reticulation experiment      # the relaxed-clock refutation
./target/release/reticulation align-error     # the aligner's admissibility ladder
curl -L -o test.parquet \
  https://huggingface.co/datasets/hiyata/Virus-Host-Genomes/resolve/main/test.parquet
python scripts/reticulata_hf_populations.py --parquet test.parquet \
  --out data/reticulata_phase3/populations --taxa 16
./target/release/reticulation real-data \
  --populations ../../data/reticulata_phase3/populations \
  --emit-lean ../../data/reticulata_phase3/emitted
\end{lstlisting}

The first reproduces Numerics~12.1 and~12.2; add \texttt{-{}-generations 5
-{}-seeds 3 -{}-max-quadruples 60000} for the thirty-two-taxon table, which
searches $35\,960$ quadruples per dataset and refuses to subsample. The second
reproduces Numerics~12.3 in about two minutes and \emph{prints the admissibility
threshold it measures}; the third and fourth reproduce Numerics~12.4. All are
deterministic given their seeds, and the population extraction is deterministic
given the parquet, whose SHA-256 is recorded in
\file{data/reticulata\_phase3/populations/provenance.json}.

\texttt{real-data} applies the threshold \texttt{align-error} measured, as
defaults to \texttt{-{}-max-gap-fraction} and \texttt{-{}-max-insertion-rate}.
Raising either is a deliberate act with a known cost: at a heavier load the
aligner alone manufactures four-point violations on a genealogy that is exactly
a tree, and the command would then be reporting that instead of the population.

\subsection*{3. Check the Lean proofs}

\begin{lstlisting}[language=bash,style=lean,morekeywords={python}]
python -m anse.formal.lean_runner formal/ANSE/ProofreadingStageRemoval.lean
python -m anse.formal.lean_runner formal/ANSE/GeneticsTextbook.lean
python -m anse.formal.lean_runner formal/ANSE/InteractingStagesAndTransport.lean
\end{lstlisting}

Each exits~$0$ only if every declared theorem carries a \lean{\#print axioms}
line, no proof depends on \lean{sorryAx}, and no proof uses an axiom outside
$\{$\lean{propext}, \lean{Classical.choice}, \lean{Quot.sound}$\}$. Compiling is
not sufficient: a file full of \texttt{sorry} compiles perfectly well.

\begin{table}[htbp]
\centering
\caption{Lean modules and their gate status at the time of writing.}
\label{tab:leanstatus}
\begin{tabularx}{\linewidth}{@{}l r r X@{}}
\toprule
\textbf{Module} & \textbf{Decls} & \textbf{Clean} & \textbf{Used in} \\
\midrule
\lean{ANSE.ProofreadingStageRemoval} & 7 & 7 & §\ref{sec:nogo}, §\ref{sec:refutation} \\
\lean{ANSE.GeneticsTextbook}        & 11 & 11 & §\ref{sec:proofreading}, §\ref{sec:mmr}, §\ref{sec:dnatop} \\
\lean{ANSE.InteractingStagesAndTransport} & 11 & 11 & §\ref{sec:gaps} \\
\bottomrule
\end{tabularx}
\end{table}

\subsection*{4. Build the book}

\begin{lstlisting}[language=bash,style=lean,morekeywords={python}]
python textbook/build.py                 # the whole book
python textbook/build.py --section 4     # one section, alone
python textbook/build.py --all-sections  # each section separately
python textbook/build.py --list          # inventory
\end{lstlisting}

The book is one file per section. Building a single section reuses the master
preamble verbatim, so a section that compiles alone also compiles in the book;
that is what makes section-at-a-time review trustworthy. The build fails on any
LaTeX error, undefined reference, undefined citation or overfull box. Solo builds
relax the reference check, since a section cannot resolve links into its
neighbours.

\subsection*{What is not reproducible here}

Stated plainly, because an appendix like this invites the impression that
everything has been checked:

\begin{itemize}[leftmargin=*,itemsep=2pt]
  \item \textbf{The literature rates.} The measurements of
        Section~\ref{sec:refutation} are read from published tables. We have not
        repeated the experiments, and nothing in this repository could.
  \item \textbf{The yeast-to-human correspondence, functionally.}
        \S\ref{sub:ortho} now verifies that \texttt{P301} is the orthologous
        \emph{position} of \texttt{P286}. That the two enzymes respond
        equivalently to the substitution is still assumed, and no alignment can
        settle it.
  \item \textbf{The synthesis ratio.} \lean{transport_needs_synthesis_bound}
        names the hypothesis Section~\ref{sec:mmr} needs --- that the variant
        synthesises no more of the genome than the exonuclease-null enzyme ---
        and nothing in this repository measures it. It is the most useful
        missing experiment in the book.
  \item \textbf{Theorem~\ref{thm:cwf}} (Călugăreanu--White--Fuller) is cited, not
        proved, and is not formalised.
  \item \textbf{Peer review.} This book has \emph{not} been peer reviewed.
        An earlier generation of this material carried a description of a
        ``five-round'' review that was generated by a script making no model
        calls; that description was withdrawn, and nothing has replaced it.
        The internal review described in §\ref{app:review} is not peer review
        and is not offered as a substitute: it was performed by the same
        authorship that wrote the book, which is precisely the independence that
        peer review exists to supply.
\end{itemize}

\subsection*{The internal review of this edition}\label{app:review}

Before release, every section was read against the data files and the primary
sources, each quoted number recomputed, and every author--year citation checked
against a primary record. This is recorded because the review changed the book's
content rather than merely confirming it. What it found:

\begin{enumerate}[leftmargin=*,itemsep=3pt]
  \item \textbf{\texttt{V411L} is not a neutral variant}
        (§\ref{sec:refutation} control~2, Worked Example~7.1). The first edition
        called it functionally neutral on the strength of one MMR-proficient
        yeast measurement. Barbari \emph{et al.}\ (2022) report a robust mutator
        phenotype for the analogous allele once mismatch repair is removed. The
        claim was corrected and the reasoning error named.
  \item \textbf{Four unsourced ratios.} The variant panel of
        Figure~\ref{fig:panel} carried hard-coded values ($21$, $12.5$, $9.1$,
        $0.55$) that could not be reproduced from
        \file{data/pole_proofreading_literature.json} by any normalisation. They
        are now computed from that file ($10.1$, $5.6$, $4.3$, $0.42$, with
        \texttt{pol2-D290V} added), with the yeast allele names corrected ---
        \texttt{pol2-F382S} and \texttt{pol2-S474F}, not the human
        \texttt{F367S} and \texttt{S459F} --- and whether an allele breaks the
        ceiling is now decided by $\rho > 1$ rather than asserted alongside it.
  \item \textbf{A wild-type rate wrong by a factor of ten}
        (Table~\ref{tab:kane}). Printed as $1.5\times10^{-8}$ against the
        provenance record's $15$. Corrected; $\rho = 2300/44$ and therefore
        Theorem~\ref{thm:refuted} were unaffected.
  \item \textbf{An unstable statistic reported as a result}
        (§\ref{sec:accumulation}). The variance-to-mean ratio of $346.7$ moves
        by seventy-fold under reseeding and sits far below the closed form. The
        section now reports the spread and the closed form, and
        §\ref{sub:trap} explains why the sample variance is systematically
        biased downward --- which is a sharper version of the lesson the section
        was written to teach.
  \item \textbf{An entailment that is conditional}
        (§\ref{sec:mmr}, Remark~\ref{rem:transport}). Theorem~\ref{thm:forced}
        is sound, but applying it transports an in vitro inequality into a cell,
        and Xing \emph{et al.}'s hyperactivity result gives a concrete reason that
        step may fail. The capacity drop is now presented as conditionally
        forced, with the condition in the ledger as an assumption.
  \item \textbf{No reference list.} Sources were cited author--year in prose
        with no list anywhere. §\ref{app:refs} supplies one, with DOIs verified
        against primary records and three entries explicitly marked as carrying
        no verified DOI.
  \item \textbf{A factor of two in a model statement}
        (§\ref{sec:accumulation}). Assumption~\ref{model:ld} said $\mu$ per
        division where the simulation used $2\mu$; the text now states $\nu =
        2\mu$ and the closed forms carry it.
\end{enumerate}

Items 1, 2, 3 and 7 are errors of fact. Items 4, 5 and 6 are overclaiming of the
kind this book is organised against, found in the book itself. Readers should
assume the list is incomplete.

\subsection*{The third edition's own review}

The same three checks were applied again when Section~\ref{sec:gaps} was added,
and found two more things --- both in the new material, which is the expected
place for them.

\begin{enumerate}[leftmargin=*,itemsep=3pt]
  \item \textbf{A frame-sampling artefact in our own experiment}
        (\S\ref{sub:strideartefact}). The shuffle null subsamples the chaos-game
        cloud, and the first run chose a stride of $6$ for \texttt{POLD1}, which
        shares a factor with the codon period. Worked Example~8.1 of this book
        warns about exactly that, in those words. The run produced an apparently
        significant \texttt{POLD1} signal which was an artefact of our own
        sampling. The nested nulls diagnosed it, the stride is now forced
        coprime with~$3$, and both the script and
        \file{tests/test_tda_null_and_orthologue.py} assert it.
  \item \textbf{The wrong statistic reported as the evidence}
        (Remark~\ref{rem:diagonal}). ``Human $286$ aligns to yeast $301$'' is
        weak, because a global alignment of similar-length sequences follows the
        diagonal: the composition-matched shuffle control lands on $302$, one
        residue away. The script now gates on local conservation instead, which
        has no such prior. Had we reported the position, the control would have
        appeared to pass while establishing nothing.
\end{enumerate}

Outside the textbook, the repository's \file{README.md} carried a badge reading
\texttt{Claude\_Opus\_4.6\_Peer\_Review — 5/5 Rounds PASSED | Score 1.00}, a
diagram node repeating it, and a table citing the review files as a credential.
That review made no model call: its referee findings were hard-coded strings and
its scores were pass-fractions of local string checks. It had been withdrawn
from the deposited papers in a previous edition but not from the front page. The
badge now reads \texttt{NOT peer reviewed}.
\file{docs/RESEARCH_DIRECTIONS_AND_ROADMAP.md} also still presented the refuted
stage-removal mechanism as an established result in four places; each is now
marked inline, with the arithmetic left visible and the model labelled false.

\subsection*{Environment}

Figures and data require Python with \texttt{numpy}, \texttt{scipy},
\texttt{matplotlib} and \texttt{gudhi} ($\ge 3.13$). Lean requires the pinned
toolchain of the enclosing project; the gate script locates it. The book requires
a \TeX\ distribution with \texttt{tcolorbox}, \texttt{listings},
\texttt{tabularx} and \texttt{siunitx}.
